\subsection{Rejection sampling acceptance rates}
\label{sec:rejection-sampling}

To probe Viterbo's conjecture computationally, we need large datasets of
random convex polytopes in~$\mathbb{R}^4$. We generate these by rejection
sampling: sample a candidate halfspace representation, then validate that it
defines a bounded, irredundant polytope.

\subsubsection*{Sampling procedure}

A polytope~$K \subset \mathbb{R}^4$ is specified in halfspace representation
\[
  K = \bigl\{\, x \in \mathbb{R}^4 \;\big|\; n_i \cdot x \le h_i,\;
  i = 1, \ldots, F \,\bigr\},
\]
where~$n_i \in S^3$ are outward unit normals and~$h_i > 0$ are heights
(positive heights ensure the origin lies in the interior).

For a given facet count~$F$ and height range~$[h_{\min}, h_{\max}]$, each
sample attempt proceeds as follows:
\begin{enumerate}
  \item Sample~$F$ independent unit normals uniformly on~$S^3$
    (via normalized standard Gaussian vectors in~$\mathbb{R}^4$).
  \item Sample~$F$ independent heights uniformly from~$[h_{\min}, h_{\max}]$.
  \item Validate the candidate~$(n_1, \ldots, n_F, h_1, \ldots, h_F)$:
  \begin{enumerate}
    \item \emph{No duplicate normals}: reject if~$\|n_i - n_j\| < \varepsilon$
      for any~$i \ne j$.
    \item \emph{Boundedness}: for every triple~$(i,j,k)$ of normals, compute
      the 1D kernel~$d$ of the $3 \times 4$ matrix~$[n_i;\, n_j;\, n_k]$
      via the generalized cross product. Verify that some normal outside the
      triple has positive inner product with both~$d$ and~$-d$.
    \item \emph{Vertex enumeration}: solve all~$\binom{F}{4}$ systems of
      $4$ active constraints, keep feasible solutions, deduplicate.
    \item \emph{Irredundancy}: for each facet, check that the incident
      vertices have affine rank~$3$ (via SVD).
  \end{enumerate}
  \item If all checks pass, accept the polytope.
\end{enumerate}

\subsubsection*{Acceptance rate sweep}

We measured acceptance rates across a grid of parameters: facet
counts~$F \in \{5, 6, \ldots, 10\}$ and three height
ranges~$[h_{\min}, h_{\max}]$. For each configuration, 1000~attempts were
made with a fixed random seed.

\begin{table}[ht]
\centering
\begin{tabular}{r ccc}
\toprule
& \multicolumn{3}{c}{Acceptance rate} \\
\cmidrule(lr){2-4}
$F$ & $h \in [0.5, 2.0]$ & $h \in [0.1, 5.0]$ & $h \in [0.8, 1.2]$ \\
\midrule
 5 & 0.057 & 0.057 & 0.057 \\
 6 & 0.168 & 0.162 & 0.170 \\
 7 & 0.334 & 0.305 & 0.344 \\
 8 & 0.450 & 0.371 & 0.506 \\
 9 & 0.489 & 0.332 & 0.602 \\
10 & 0.504 & 0.283 & 0.729 \\
\bottomrule
\end{tabular}
\caption{Acceptance rates for rejection sampling of random 4-polytopes
  with~$F$ facets. Each cell reports the fraction of 1000 attempts that passed
  all validation checks.}
\label{tab:acceptance-rates}
\end{table}

\subsubsection*{Observations}

\paragraph{Acceptance increases with~$F$ for narrow height ranges.}
For~$h \in [0.8, 1.2]$, the rate climbs from~5.7\% at~$F = 5$ to~72.9\%
at~$F = 10$. More facets provide more directions to ``close off'' the
polytope, making boundedness easier to satisfy.

\paragraph{Acceptance is insensitive to height range at low~$F$.}
At~$F = 5$, all three height ranges give exactly the same rate (5.7\%).
The dominant rejection cause at low facet count is the boundedness check,
which depends only on the normal directions, not the heights.

\paragraph{Wide height ranges hurt at high~$F$.}
At~$F = 10$, the rate drops from~72.9\% for~$h \in [0.8, 1.2]$ to~28.3\%
for~$h \in [0.1, 5.0]$. Once boundedness is likely satisfied, the irredundancy
check becomes the bottleneck: extreme height ratios produce degenerate vertex
configurations that fail the affine rank test.

\paragraph{Practical consequence.}
For building datasets with~$F \le 8$, rejection sampling is efficient enough
(acceptance~$\ge 37\%$) without biasing the distribution. For~$F = 5$, the
5.7\% rate means roughly 18 attempts per accepted polytope --- still practical
for datasets of a few thousand samples.

{\color{red}%
\textbf{Stub data.}
Timing data in this table uses stub implementations (constant 1.0 for volume
and capacity, busy-spin for ${\sim}1$\,ms). The acceptance rates and
acceptance/rejection counts are real. The systolic ratio experiment
(Section~\ref{sec:systolic-ratios}, when written) will use real
implementations from the \texttt{hk2017} crate.}
